% ===== §17 =====
\section{Intimate Relation as a Special Site of Cultural Emergence}
\label{sec:special}

\noindent
The first part established what culture is and how it emerges, in general. This part turns to the site for which the general account was built and from which this series takes its name: the intimate relation between two people. The turn is not a change of subject but a descent from the universal to the particular that both instances it and exceeds it, and the descent is the heart of the paper. For if the general ontology of the first part is right, then the intimate dyad is not merely one instance of cultural emergence among others, a small example of a large phenomenon; it is the \emph{minimal and most exposed} instance, the smallest system in which a culture can arise at all, and therefore the site at which the mechanism of emergence, buried elsewhere under institutions and centuries, can be watched almost as it happens. The disproportion the introduction noted---that a philosophy of intimacy should have discussed culture in general at such length and the intimate case so little---is corrected here, and corrected deliberately, for the general was the long approach and the particular is the destination.

\subsection*{Why the dyad is a special site, not merely a small one}

It would be possible to treat the intimate dyad as simply a culture in miniature, subject to the same account as any other culture only at a smaller scale, and to read off its features by shrinking the features of macro-culture. This is exactly what the paper does \emph{not} do, and the refusal is the methodological hinge of the whole part. The dyad is not a small culture; it is a \emph{special} one, possessed of structural features that belong to it alone and appear in no larger cultural formation, features that do not scale up and cannot be derived by shrinking the large. The reason was given, in outline, at the close of the definition: the mechanism of emergence may be stated generally, but the specific conditions under which it operates in the dyad---above all the absence of any external authority to fix meaning, and the two-person structure itself---have no counterpart at the scale of a people, which possesses a full apparatus of traditions and institutions where the couple possesses nothing but the two of them. The dyad is therefore special in the strict sense: not a reduced version of the general case, but a case with its own irreducible mechanisms, which the general account illuminates but does not exhaust. To describe those mechanisms is the work of this part, and it is where the paper means to make its distinctive contribution, for these mechanisms have not, to the author's knowledge, been gathered and named as the mechanisms of a culture.

\subsection*{The five pillars}

Five such features organize the phenomenology that follows, and they are named here in advance so that their treatment is read as the articulation of a single structure rather than a list of observations. The first is \emph{the dyad and the vacant big Other}: the intimate relation is the minimal cultural unit, two people with no external third---no tradition, no institution, no majority---to fix the meaning of what they generate, so that the two must serve as each other's symbolic authority, which is at once what makes their private culture possible and what makes it fragile. The second is \emph{the double interweaving of drive}: the two structures of desire are not parallel but mutually implicated, each entering the other's fantasy, so that the culture they generate is the sediment of an interweaving that exists nowhere outside the relation. The third is \emph{high-density sedimentation}: the density of shared living---daily, bodily, high in frequency and long in duration---lets intimate culture sediment at a speed no other cultural formation approaches, so that one can almost watch a private form come into being. The fourth is \emph{the special configuration of the three registers}: in intimate relation the imaginary, the symbolic, and the real are configured with a closeness and intensity found in no other cultural site, so that the three-register account of the first part reaches here its fullest embodiment. The fifth is \emph{fragility and reversibility}: an intimate culture can collapse entire, in a way a macro-culture cannot, because its whole existence depends on the continued participation of exactly two people, and this fragility is the price of the very features that make it possible.

\subsection*{Method: a return to the phenomena}

One methodological commitment governs this part, and it should be stated before the description begins. The first part was theoretical; it argued, mapped, defined. This part suspends that posture, as far as it can, and returns to the phenomena themselves---to intimate culture as it is actually lived, experienced, and understood, before it is theorized. This is the phenomenological \emph{epoch\'e} in the modest sense the paper can honor: a bracketing of the theoretical apparatus so that the mechanisms may be allowed to show themselves from the phenomena rather than be deduced from the framework. The framework of the first part is not abandoned---it is what lets the phenomena be understood once they have shown themselves---but it is held back at the moment of description, so that the account earns its mechanisms from the lived material and does not merely stamp the general theory onto a smaller case. The order of the part follows this commitment: first the epistemology of intimate culture, how such culture is known and by whom; then its mechanisms, the five pillars in their working; then its hermeneutics, how the meaning of intimate culture is interpreted, and how, at its edge, it is drawn back into negotiation with the outer symbolic order, which returns the particular to the general and opens the way to the ethics and political economy of the final part. The description begins with the phenomena as they are given, and asks first not what they are but how they are known.
